\section{Experiments}

\subsection{Data}

We evaluate on four full-route city bundles.
Singapore uses LTA DataMall passenger-volume and OD bus data with traffic speed bands.
Austin/CapMetro uses GTFS schedules with demand proxied from scheduled headways.
Halifax Transit uses route-level APC summaries apportioned over GTFS route-direction patterns.
MBTA uses Fall 2025 stop-level boardings and alightings apportioned over GTFS route-direction patterns.
\Cref{tab:dataset} summarizes the bundles and their evidence levels.
Because these sources are heterogeneous, \Cref{tab:target-audit-main} reports the target-construction audit in the main text.
The table separates observed evidence from schedule-derived proxies and reports route-level target residual quantiles, making explicit that the primary target is static-derived rather than direct counterfactual operational ground truth.
\Cref{fig:validation-protocol} summarizes the data-construction and validation boundary: source residual labels are used for fitting, while target transition labels are blocked until the final reporting step.

\begin{figure}[!t]
    \centering
    \includegraphics[width=\textwidth]{validation_protocol_schematic.png}
    \caption{Validation protocol and data-construction boundary. Open agency data are converted into full-route city bundles, paired simulator outputs $f_0(s,a)$ and static-derived targets $f_c(s,a)$, and source-city residual labels. In strict leave-one-city-out validation, target covariates and unlabeled summaries may be used for source weighting and feature-set selection, but held-out target transition labels are used only after selection is frozen for final reporting.}
    \label{fig:validation-protocol}
\end{figure}

\begin{table}[t]
\centering
\scriptsize
\resizebox{\textwidth}{!}{\begin{tabular}{lllllllllllll}
\toprule
city & env\_key & lines & stops & segments & timetables & transitions & median\_route\_km & median\_speed\_mps & median\_headway\_s & peak\_demand\_hour & demand\_source & traffic\_source \\
\midrule
Singapore & singapore\_lta\_all & 789 & 25873 & 25906 & 56545 & 1243488 & 17.0 & 9.58 & 690.0 & 7 & LTA DataMall PV/ODBus 202604, WEEKDAY monthly totals divided to typical weekday & LTA TrafficSpeedBands snapshot nearest-matched to stop segments \\
Austin / CapMetro & austin\_capmetro\_all & 234 & 6464 & 6464 & 28329 & 310272 & 16.1678 & 6.19 & 1200.0 & 14 & GTFS schedule headway proxy OD; no public no-key stop-level ridership found & GTFS schedule-derived segment speed \\
Halifax Transit & halifax\_transit\_all & 176 & 5551 & 5551 & 8710 & 266448 & 12.3693 & 6.84 & 1140.0 & 15 & Halifax route-level half-hour APC apportioned across GTFS route-direction-patterns & GTFS schedule-derived segment speed \\
MBTA & mbta\_all & 940 & 20195 & 20195 & 67091 & 969360 & 7.6405 & 4.01 & 600.0 & 7 & MBTA Fall 2025 stop-level boardings/alightings apportioned across GTFS route-direction-patterns & GTFS schedule-derived segment speed \\
\bottomrule
\end{tabular}
}
\caption{Open transit city bundles. Evidence levels distinguish observed passenger/traffic data from schedule-derived proxies.}
\label{tab:dataset}
\end{table}

\begin{table}[t]
\centering
\scriptsize
\resizebox{\textwidth}{!}{\begin{tabular}{lllllllll}
\toprule
city & demand\_evidence & traffic\_evidence & lines & transitions & uncal\_mse & line\_mse\_p25 & line\_mse\_p50 & line\_mse\_p75 \\
\midrule
Singapore & observed stop OD/PV monthly & observed speed bands & 789 & 1243488 & 2256.5339 & 220.173 & 324.7508 & 541.4336 \\
Austin / CapMetro & schedule proxy & schedule derived & 234 & 310272 & 1301.3027 & 242.2482 & 615.8957 & 3043.8061 \\
Halifax Transit & route APC apportioned & schedule derived & 176 & 266448 & 1111.9223 & 199.2164 & 620.2908 & 1614.3016 \\
MBTA & stop board/alight apportioned & schedule derived & 940 & 969360 & 1216.3716 & 180.9625 & 459.6093 & 1040.0624 \\
\bottomrule
\end{tabular}
}
\caption{Target-construction audit. The evidence columns identify observed, apportioned, and schedule-derived components; route-level residual quantiles show that simulator-target gaps are heterogeneous rather than a single constant correction.}
\label{tab:target-audit-main}
\end{table}

\subsection{Validation protocols}

\paragraph{Primary dynamics validation.}
The main dynamics experiment is strict leave-one-city-out.
Each city is held out once as the target, and all remaining cities are used as sources.
The metric is total MSE averaged over the nine transition outputs.
Ratios below 1 indicate improvement over the dense \Htwo-style residual baseline.

\paragraph{Single-source zero-shot diagnostic.}
We also run every ordered one-source-to-one-target city pair.
Only the source city provides residual labels; the target city provides the same network and static summary inputs used by the target-label-free selector, but no target transition labels before reporting.
This diagnostic is intentionally harder and is not expected to be uniformly positive.
Its role is to establish the minimum-information baseline before adding source diversity or target passive history.

\paragraph{Target offline adaptation / few-shot curve.}
The target-route budget sweep is interpreted as offline adaptation, not zero-shot transfer.
For each held-out target city, $0\%$, $1\%$, $5\%$, $10\%$, $25\%$, or $100\%$ of target route groups are added to the source sufficient statistics, and evaluation is performed on the remaining target route groups except at the $100\%$ in-sample oracle reference.
This approximates the practical deployment setting in which a new city has a small amount of passive AVL/APC or schedule-derived historical labels before control deployment.
We additionally run same-information few-shot adaptation baselines motivated by AdaRL and FANsRL.
The AdaRL-style baseline uses the target budget to select which mechanism groups should be adapted and then refits only those groups.
The FANsRL-style baseline keeps the transferred mechanism parameters fixed and estimates a low-dimensional per-output target context bias.
These are not direct executions of the original CartPole/Atari or MuJoCo codebases; they are bus-residual evaluators that use the same target labels and the same source-city statistics as \ours.

\paragraph{One-step policy proxy.}
The one-step policy experiment uses the same source-target splits and evaluates action selection under headway-stress states.
We report reward/regret and traffic-domain metrics.
Bunching rate is the fraction of decisions in which either adjacent headway is below $0.5$ times the target headway.
Large-gap rate is the fraction in which either adjacent headway exceeds $1.5$ times the target headway.

\paragraph{Reviewer-facing robustness checks.}
We include diagnostics targeted at common review objections.
First, a target-construction audit separates observed, apportioned, schedule-proxy, and deterministic target components.
Second, the target-route offline-adaptation sweep adds $0\%$, $1\%$, $5\%$, $10\%$, $25\%$, or $100\%$ of held-out target-city routes to the source sufficient statistics and evaluates on the remaining target routes; the $100\%$ point is an in-sample oracle reference.
Third, per-mechanism errors decompose total MSE into headway/gap, demand/dwell/reward, and speed groups.
Fourth, mechanism ablations compare the proposed parent sets against wrong mechanism groupings, action/time-only features, and dense matched-capacity variants.
Fifth, source-weight sensitivity varies $T \in \{0.5,1.0,2.0\}$ and $\lambda \in \{0,0.05,0.1\}$, while source-subset and pairwise-source checks remove major source families.
Sixth, generator robustness perturbs output groups or adds measurement-noise terms to the evaluation sufficient statistics.

\paragraph{Full-module guarded-ensemble validation.}
The main paper-suite tables use deterministic sufficient-statistics variants for reproducibility.
We additionally report a full-module diagnostic that adds rigid-first residual deltas, deterministic static-local features, source-city leave-one validation, and the guarded source-city ensemble.
For the ensemble weights, a source-only feature-set selector chooses among simulator-output-only, observation-only, observation-plus-simulator, local-only, full-summary, and uniform-weight alternatives before target evaluation.
We also add a same-information dense \Htwo baseline that uses the same unlabeled target summary to weight source-city dense residual experts.
This diagnostic uses six configured cross-city splits: the four strict leave-one-city-out targets plus two source-family transfer splits used to stress the source selector.
Because this diagnostic contains repeated target cities, it is interpreted as a selector stress test rather than as six independent city-level targets.

\paragraph{External real APC/AVL validation.}
To partially address the static-derived target limitation, we add an external Austin/CapMetro validation slice from public raw APC stop-event data.
The July--December 2021 and August 2016 files contain observed stop events with boardings, alightings, vehicle load, dwell time, vehicle position, route, trip, stop, and vehicle identifiers.
We filter records with valid stop, trip, GPS, and import-quality flags, reconstruct stop-level headways by route-direction-stop event times, and evaluate passive action-$0$ one-step dynamics against the observed next stop event.
The no-correction baseline in this experiment is a state-informed passive transition using the APC-derived current state, not a schedule-only free-running simulator.
This is real historical APC/AVL validation, but it is still offline passive validation rather than a field control deployment because counterfactual hold actions are not observed.
In observational AVL/APC logs, a counterfactual hold outcome would require knowing what would have happened in the same state under a hold action that was not actually taken; such labels require randomized or quasi-randomized field interventions, an experimentally validated simulator, or another credible causal design.
We also run a few-shot APC adaptation sweep at $0\%$, $1\%$, $5\%$, $10\%$, and $25\%$ of Austin line groups.
This sweep compares source-only transfer, source-plus-target residual refitting, target-only residual fitting, a residual gate that learns how much of the source correction to trust, and a target residual-bias adapter.
The gate and bias-adapter variants are lightweight AdaRL/FANsRL-style checks of what changes should be adapted; they are not full reimplementations of those neural RL systems.
Because route-random budgets are not the most natural deployment unit, we also add a chronological target-history sweep using the first $0$, $1$, $6$, $24$, or $168$ hours of Austin APC history for adaptation and all later APC transitions for evaluation.

\paragraph{Auxiliary executable rollout.}
Sampled BusSimEnv rollout is reported as an auxiliary sanity check on runnable line-level environments generated from each city bundle.
For this rollout, each policy is evaluated on the same line-level random seed, including Python signal-delay randomness, and live observations are mapped back to the static-transition feature schema before model scoring.
The policy score uses the BusSimEnv linear headway reward rather than the static one-step reward.

\paragraph{Full-line libsumo holding diagnostic.}
We also generate executable SUMO scenarios with explicit bus-stop dwell points for every validated line in each city bundle.
The diagnostic uses a four-hour window with at most eight departures per line and evaluates no-hold, dense \Htwo-MPC, weighted \ours-MPC, threshold equalization, retuned Daganzo, and a target-summary gated Daganzo prior.
The target-summary prior uses only static target route and schedule summaries to choose whether a conservative or aggressive Daganzo gain is appropriate; it does not use target transition labels or the held-out closed-loop result for fitting.
This rollout is a counterfactual simulator diagnostic, not field-calibrated evidence, but it tests whether the proposed mechanism-state idea should control a safe prior or a free-form residual action.

\paragraph{Native-protocol RL cross-checks.}
Finally, we run official RESCO and patched LibSignal RL baselines under each library's native state, reward, action, and logging definitions.
The official RESCO table trains \texttt{IDQN} and upstream-supported \texttt{MPLight} profiles for 100 episodes and evaluates five test episodes on the extended RESCO scenario set.
The LibSignal table trains \texttt{DQN}, \texttt{PressLight}, \texttt{FRAP}, and \texttt{MPLight} for 100 episodes on \texttt{sumo1x1}, \texttt{sumo1x3}, and \texttt{sumo4x4}, with a final 3600-step test.
These rows answer a reviewer-facing reproducibility question: standard RL runners are connected and run at full budget.
They are not same-protocol comparisons against \ours, because their observation spaces, rewards, action timing, and reported metrics differ from the phase-MPC stress-test protocol.

\begin{table}[t]
\centering
\small
\resizebox{\textwidth}{!}{\begin{tabular}{llll}
\toprule
experiment & paper\_role & unit & scope\_note \\
\midrule
strict\_leave\_one\_city\_out & primary cross-city dynamics evidence & city target & four unique targets; no duplicated target counted as independent \\
generator\_robustness & bias/noise defense & scenario x city target & static-derived target perturbation; not real AVL/APC trajectory ground truth \\
source\_subset\_robustness & data-source confound defense & city subset & reports excluding Singapore and excluding Austin proxy-demand subsets \\
cross\_city\_policy\_validation & primary one-step policy evidence & city target & one-step lookahead, not live SUMO rollout \\
sampled\_rollout & auxiliary sanity check only & BusSimEnv episode & 10 policy episodes across runnable line envs; do not use as primary claim \\
\bottomrule
\end{tabular}
}
\caption{Experimental scope and paper role. The primary evidence is strict city-level leave-one-out dynamics; one-step policy and sampled rollout results are secondary control-facing diagnostics.}
\label{tab:scope}
\end{table}
